\documentclass[10pt,a4paper]{article}
\usepackage[a4paper,margin=14mm,top=12mm,bottom=14mm]{geometry}
\usepackage{fontspec}
\setmainfont{Liberation Sans}
\usepackage{polyglossia}
\setdefaultlanguage{french}
\usepackage{amsmath}
\usepackage{xcolor}
\usepackage{graphicx}
\usepackage[most]{tcolorbox}
\usepackage{tabularx,booktabs,array,multirow}
\usepackage{enumitem}
\usepackage{fancyhdr}
\usepackage{lastpage}
\usepackage[hidelinks]{hyperref}
\definecolor{uimm}{HTML}{005677}
\definecolor{cyan}{HTML}{13A2DD}
\definecolor{fond}{HTML}{E6F2FA}
\definecolor{ok}{HTML}{2E7D32}
\definecolor{ko}{HTML}{B3261E}
\definecolor{violet}{HTML}{7B1FA2}
\setlength{\parindent}{0pt}
\setlength{\parskip}{3pt}
\setlist{nosep}
\setlist[itemize]{leftmargin=13pt,label=\textcolor{cyan}{\textbullet}}
\newcommand{\titre}[1]{\vspace{5pt}{\color{uimm}\large\bfseries #1}\par\vspace{-3pt}{\color{cyan}\rule{\linewidth}{1.2pt}}\par\vspace{0pt}}
\tcbset{boite/.style={colback=fond,colframe=uimm,boxrule=0.6pt,arc=2pt,left=6pt,right=6pt,top=4pt,bottom=4pt},
        defi/.style={colback=violet!6,colframe=violet,boxrule=0.8pt,arc=2pt,left=6pt,right=6pt,top=4pt,bottom=4pt,
                     title={\bfseries Défi},coltitle=white,colbacktitle=violet,fonttitle=\small}}
\newcommand{\bandeau}[3]{%
\begin{tcolorbox}[colback=uimm,colframe=uimm,arc=2pt,left=4pt,right=8pt,top=4pt,bottom=4pt]
\begin{minipage}[c]{18mm}\setlength{\fboxsep}{1.2pt}\fcolorbox{white}{white}{\includegraphics[height=16.5mm]{logo_uimm.jpg}}\end{minipage}\hspace{4mm}%
\begin{minipage}[c]{\dimexpr\linewidth-22mm\relax}\color{white}
{\small #1}\\[1pt]{\LARGE\bfseries #2}\\[2pt]{\large #3}\hfill{\small CQPM Technicien qualité}
\end{minipage}
\end{tcolorbox}}
\newcommand{\reponse}[1][1]{\par\vspace{2pt}{\color{gray!60}\foreach \i in {1,...,#1}{\rule{\linewidth}{0.3pt}\par\vspace{9pt}}}}
\usepackage{pgffor}
\pagestyle{fancy}\fancyhf{}
\renewcommand{\headrulewidth}{0pt}
\newcommand{\pied}[1]{\fancyfoot[L]{\footnotesize\color{uimm}#1}\fancyfoot[R]{\footnotesize\color{uimm}Page \thepage{} sur \pageref{LastPage}}}

\usepackage{tikz}
\usepackage{longtable}
\renewcommand{\bandeau}[3]{%
\begin{tcolorbox}[colback=uimm,colframe=uimm,arc=2pt,left=4pt,right=8pt,top=4pt,bottom=4pt]
\begin{minipage}[c]{18mm}\setlength{\fboxsep}{1.2pt}\fcolorbox{white}{white}{\includegraphics[height=16.5mm]{logo_uimm.jpg}}\end{minipage}\hspace{4mm}%
\begin{minipage}[c]{\dimexpr\linewidth-22mm\relax}\color{white}
{\small #1}\\[1pt]{\LARGE\bfseries #2}\\[2pt]{\large #3}
\end{minipage}
\end{tcolorbox}}
% en-tête d'une fiche : code, titre, question, bloc, durée, organisation, niveau
\newcommand{\entete}[7]{%
\begin{tcolorbox}[colback=uimm,colframe=uimm,arc=2pt,left=4pt,right=6pt,top=4pt,bottom=4pt]
\begin{minipage}[c]{18mm}\setlength{\fboxsep}{1.2pt}\fcolorbox{white}{white}{\includegraphics[height=16.5mm]{logo_uimm.jpg}}\end{minipage}\hspace{3mm}%
\begin{minipage}[c]{\dimexpr\linewidth-58mm\relax}\color{white}\raggedright
{\small Base SPC · #4}\\[1pt]{\Large\bfseries #1 · #2}\\[2pt]{\itshape #3}
\end{minipage}\hfill
\begin{minipage}[c]{33mm}\color{white}\small\raggedleft
Durée : \textbf{#5}\\ \textbf{#6}\\ Niveau #7 sur 3
\end{minipage}
\end{tcolorbox}}
\newcommand{\tagE}{{\setlength{\fboxsep}{1.5pt}\colorbox{cyan}{\color{white}\scriptsize\bfseries Ellistat}}}
\newcommand{\tagW}{{\setlength{\fboxsep}{1.5pt}\colorbox{uimm}{\color{white}\scriptsize\bfseries Page web}}}
\tcbset{pred/.style={colback=orange!6,colframe=orange!80!black,boxrule=0.6pt,arc=2pt,left=6pt,right=6pt,top=3pt,bottom=1pt,
                    title={\bfseries Je prédis, avant de commencer},coltitle=white,colbacktitle=orange!80!black,fonttitle=\small},
        retenir/.style={colback=fond,colframe=uimm,boxrule=0.6pt,arc=2pt,left=6pt,right=6pt,top=3pt,bottom=3pt,
                    title={\bfseries À retenir},coltitle=white,colbacktitle=uimm,fonttitle=\small},
        plus/.style={colback=violet!6,colframe=violet,boxrule=0.6pt,arc=2pt,left=6pt,right=6pt,top=3pt,bottom=3pt,
                    title={\bfseries Pour aller plus loin},coltitle=white,colbacktitle=violet,fonttitle=\small},
        act/.style={colback=uimm!10,colframe=uimm,boxrule=0.5pt,arc=2pt,left=5pt,right=5pt,top=2pt,bottom=2pt,before skip=7pt,after skip=3pt},
        verif/.style={colback=orange!6,colframe=orange!80!black,boxrule=0.6pt,arc=2pt,left=6pt,right=6pt,top=3pt,bottom=3pt,
                    title={\bfseries À vérifier sur poste avant la première séance},coltitle=white,colbacktitle=orange!80!black,fonttitle=\small}}
\newcommand{\fichier}[1]{\texttt{\small #1}}
\newcommand{\menu}[1]{\textsf{\textbf{#1}}}
\newcommand{\champ}[1]{\rule{#1}{0.3pt}}

\pied{Base SPC · fiche M10 · Fabriquer son étalon}
\begin{document}
\thispagestyle{fancy}
\entete{M10}{Fabriquer son étalon}{Comment obtenir un étalon de 12,000 mm avec des cales ?}{Le moyen de mesure}{25 min}{en binôme}{2}
\begin{tcolorbox}[boite]\textbf{Le but.} Composer une valeur par adhérence de deux cales et s'en servir comme étalon au plus près de la cote contrôlée.\par\textbf{Il vous faut.} Un coffret de cales (cales de 10 et de 2 mm) · Un micromètre 0-25 mm · Chiffon non pelucheux\end{tcolorbox}
\begin{tcolorbox}[pred]Deux cales de 10 et de 2 mm collées l'une sur l'autre font-elles exactement 12,000 mm au micromètre ?\reponse\end{tcolorbox}
\titre{À faire}
\begin{enumerate}[label=\textcolor{cyan}{\textbf{\arabic*.}},leftmargin=15pt,itemsep=1.5pt]
\item Nettoyer les faces des deux cales. Les poser en croix, presser légèrement et faire glisser jusqu'à ce qu'elles collent : c'est l'\textbf{adhérence}.
\item Mesurer la pile 15 fois au micromètre, en la retirant à chaque fois. Ne pas la tenir dans la main.
\item \tagE\ \menu{Gage R\&R}, onglet \menu{Justesse}, « Ajouter une justesse » : Nombre de mesures 15 ; Valeur vraie 12,000 ; LSI 11,85 et LSS 12,00 ; k 3. Saisir les 15 mesures avec trois décimales.
\item Séparer les cales en les faisant glisser, sans tirer. Les nettoyer et les ranger.
\end{enumerate}
\vspace{3pt}\textbf{Pile 10 + 2 : mes 15 mesures (mm)}\par\vspace{2pt}
{\renewcommand{\arraystretch}{1.45}\begin{tabularx}{\linewidth}{|*{5}{>{\centering\arraybackslash}p{5mm}|>{\centering\arraybackslash}X|}}\hline
\scriptsize 1 &  & \scriptsize 2 &  & \scriptsize 3 &  & \scriptsize 4 &  & \scriptsize 5 & \\\hline
\scriptsize 6 &  & \scriptsize 7 &  & \scriptsize 8 &  & \scriptsize 9 &  & \scriptsize 10 & \\\hline
\scriptsize 11 &  & \scriptsize 12 &  & \scriptsize 13 &  & \scriptsize 14 &  & \scriptsize 15 & \\\hline
\end{tabularx}}\par
\titre{J'observe, je décide}
\begin{enumerate}[label=\textcolor{cyan}{\textbf{Q\arabic*.}},leftmargin=22pt]
\item Moyenne de la pile = \champ{18mm} mm. Écart à 12,000 = \champ{18mm} mm.\par\vspace{5pt}
\item L'étalon vaut maintenant 12,000 mm, une limite de la tolérance de la butée. Pourquoi est-ce mieux que la cale de 10 mm ?\reponse[2]
\item Comparez Cg et Cgk avec ceux obtenus sur la cale de 10 mm (M3).\reponse[1]
\end{enumerate}
\begin{tcolorbox}[retenir]On vérifie un instrument là où il sert : avec un étalon proche de la cote à contrôler.\end{tcolorbox}
\begin{tcolorbox}[plus]Composez 11,850 mm (par exemple 10 + 1,35 + 0,5) et refaites l'étude : l'instrument est-il aussi juste aux deux limites de la tolérance ?\end{tcolorbox}
\end{document}